% =============================================================================
% 5.7 本章小结
% 草稿版本：v0.1 (2026-05-07)
% 起草动机：按导师意见 3，第 5 章实验章必须包含本章小结，且原 5.3 全文
%          总结 / 5.4 研究前景两节迁出到第 6 章，第 5 章需要一个独立的
%          收口段落把 5 套评测的核心数据与统一论证结构呈现给读者。
% 字数：约 700 字
% 风格规则：v0.2 风格（无 \textbf / \textit / 引例括号 / 双引号 / 破折号）
% 引用：暂无新增引用
% 图表：无新增（核心数据已在前 6 节表格中给出，本节仅以叙述方式收口）
% 依赖宏包：booktabs / tabularx（主稿已加载）
% =============================================================================

\subsection{本章小结}
\label{sec:eval-summary}

本章在统一的实验环境下完成了 5 份独立评测集共 214 条用例的系统化
实证。\ref{sec:exp-env} 节先说明软硬件与依赖版本，并给出 5 份评测集
的统一构建原则与可重现性工程要点；
\ref{sec:intent-eval}、\ref{sec:rag-eval}、\ref{sec:bridge-eval}、
\ref{sec:code-eval} 4 节按数据流向依次报告 4 个核心模块的能力上限；
\ref{sec:e2e-eval} 节以 30 条真实场景用例的双档对照验证整套系统的
端到端价值。

5 套评测的核心数据可一表概括。意图识别评测 35 条端到端通过率
91.4\%，多地点角色识别准确率 100\%，并暴露了多轮上下文继承的
33.3\% 瓶颈。混合检索评测 60 条 Top-1 命中率 85.0\%、MRR 0.908，
并通过权重扫描得到 $\alpha=0.9$ 的最优配比与最优权重随知识库体量
漂移的工程结论。语义桥接评测 71 条在 \texttt{rule\_\bk{}only} 与
\texttt{rule\_\bk{}plus\_\bk{}rag} 两档下 \texttt{grade\_\bk{}id} 准确率与
\texttt{citation} 真实性同时达到 100\%，相对大语言模型自由发挥基线
形成 74.6 个百分点的严格命中率差距。代码执行评测 18 条在
\texttt{llm\_\bk{}with\_\bk{}code} 档数值准确率 88.9\%，比直算高出 22.2 个
百分点，并把大语言模型在伴随状态维护与离散计数任务上的内在边界
量化到具体差距数据。端到端评测 30 条 \texttt{full} 档通过率 86.7\%，
相对 \texttt{ablated} 去 RAG 对照档提升 13.4 个百分点；权威标准
引用率从 0\% 提升到 26.7\%；\texttt{decision\_\bk{}support} 类用例
通过率从 50\% 提升到 100\%。

5 套评测在数据规模、研究问题与实证视角上互不重叠，共同构成本文
论点的完整闭环。模块层 4 套评测以中间产物视角分别量化 4 个子模块
的能力上限，集成层 1 套评测以最终用户体验视角验证整套系统的端到端
价值。两类评测之间的差距既揭示了双通路 RAG 与代码生成等核心机制
的真实贡献，也在 \texttt{citation} 真实性 100\% 衰减为端到端引用
率 26.7\% 这一指标对照上暴露出集成衰减效应。本章实证发现既支撑了
本文研究工作的有效性，也为下一章的研究局限性回顾与未来工作展望
提供了直接的数据基础。
